% =====================================================================
%  METHODOLOGY / EXPERIMENTAL SETUP  (rev. 2 — post-review)
%  Drop-in LaTeX for the McBenchface/Choreo paper (acmart).
%
%  REQUIRES: \usepackage{siunitx}; and a single definition of the
%  framework name, e.g. \newcommand{\sysname}{McBenchface}. The name is
%  currently inconsistent across the paper (prose "McBenchface", Table 2
%  and code "Choreo"): pick ONE and define \sysname once. Table 2's
%  "Choreo Pipeline" must be changed to match.
%
%  CROSS-REFS: this section cites \ref{sec:framework} (->§3 "The McBenchface
%  framework"), \ref{sec:background} (->§2), \ref{sec:intro} (->§1),
%  \ref{sec:conclusion} (->§7), and \ref{sec:results} (->the unwritten Results
%  section). The current PDF's §4/§5 are outlines with no \label{}s, so confirm
%  each \label{} exists (esp. sec:framework on the §3 heading) before submission.
%
%  Scope agreed: Devices = M2 Pro + DGX Spark (GB10) only; four
%  experiments (NoOp overhead, modularity overhead, VQA bandwidth
%  contention, Self-RAG topology). Collocation is presented as a
%  framework capability exercised inside the case studies, NOT as a
%  separate independent-workload interference experiment.
%
%  NO RESULTS are stated here. Numbers/tables/verdicts go in Results.
%
%  Placeholders marked [[FILL: ...]] need a value only the authors can
%  confirm (core counts, bandwidths, software/model/dataset versions,
%  Poisson rate, run counts, seeds).
%
%  --- OPEN DECISIONS FOR AUTHORS (raised by reviewers; reconcile before submission) ---
%  D1 (claim/evidence, HIGH): the Introduction sells "collocation-aware
%     profiling of concurrent pipelines" and "scalability" as headline
%     contributions, but this section runs neither as a first-class
%     experiment. Either (a) add the independent-workload collocation cell
%     (torchvision_mixed.yml exists) + a no-code model-size sweep +
%     an MLPerf end-to-end-vs-isolated contrast (pipeline_configs/mlperf/
%     exists), or (b) soften the Intro/Contribution-4 claims to match.
%     This file currently takes a defensible middle path (relabels
%     "scalability"->"portability"; collocation = capability) and flags
%     the remaining gap with %REVIEWER comments inline.
%  D2 (confound, MED): Self-RAG monolith(1x9B) vs decomposed(3x4B)
%     conflates model size, specialization, and call count. A size-matched
%     control arm is recommended; below it is described as a control and
%     marked [[FILL: confirm control arm run]].
% =====================================================================

\section{Experimental Methodology}
\label{sec:methodology}

This section describes how we evaluate \sysname{}: the devices we run on
(\S\ref{sec:duts}), the instrumentation and the metrics derived from it
(\S\ref{sec:metrics}), and four experiments, each stated with the question it
answers (\S\ref{sec:experiments}). Results are deferred to
Section~\ref{sec:results}. The first two experiments characterize \sysname{} as
an instrument, bounding the measurement overhead the framework itself
introduces (\S\ref{sec:exp-noop},~\S\ref{sec:exp-modularity}); the latter two
apply it to end-to-end, graph-structured measurements --- including collocation
exercised \emph{within} a workload --- that accelerator-centric, single-stream
benchmarks do not capture (\S\ref{sec:exp-bandwidth},~\S\ref{sec:exp-topology}).

% ---------------------------------------------------------------------
\subsection{Devices Under Test}
\label{sec:duts}

We evaluate on two single-node machines drawn from opposite extremes of the
unified-memory design space. Both let an accelerator and the host CPU address a
shared pool of physical memory, but they differ by roughly an order of
magnitude in capacity and by their entire software stack. This contrast lets us
test \emph{portability} across the ``mid-range gap'' that rigid
\emph{Tiny}/\emph{Edge}/\emph{Datacenter} tracks leave uncovered
(Section~\ref{sec:background}), and lets us check whether a measured systems
effect is intrinsic to a workload or specific to one platform. We do \emph{not}
treat the pair as a controlled capacity axis. Capacity differs by \num{8}$\times$
(\SI{16}{\giga\byte} vs.\ \SI{128}{\giga\byte}), but memory bandwidth by only
${\approx}\num{1.4}\times$ ([[FILL: confirm SKU --- ${\approx}\SI{200}{\giga\byte\per\second}$
for a standard M2~Pro vs.\ \SI{273}{\giga\byte\per\second}]]). Device class,
power envelope, runtime, and numerical precision all change together as well.

\begin{description}
  \item[Apple M2 Pro (edge / consumer).] A \SI{16}{\giga\byte}
    unified-memory laptop SoC
    %
    ([[FILL: CPU core split, e.g.\ 10-core 6P+4E; GPU core count, e.g.\ 16-core; memory bandwidth ${\approx}\SI{200}{\giga\byte\per\second}$]]),
    %
    in which the GPU, Neural Engine, and CPU share one physical memory and its
    bandwidth. It runs three accelerator backends over that shared memory: the
    GPU via Apple MLX/Metal (the PyTorch \texttt{mps} device --- Apple's Metal
    backend, not to be confused with NVIDIA's Multi-Process Service), the
    Neural Engine (\texttt{ane}) via Core~ML, and the CPU. Software:
    Python~3.10, MLX/\texttt{mlx-lm}, Core~ML Tools, \texttt{faiss-cpu}
    ([[FILL: macOS and package versions]]). LLM weights are 4-bit quantized to
    fit comfortably within the \SI{16}{\giga\byte} budget alongside the other
    resident stages. This machine represents the resource-constrained edge
    devices the paper argues are poorly served by existing suites.

  \item[NVIDIA DGX~Spark (large unified memory).] A GB10 Grace--Blackwell
    platform with \SI{128}{\giga\byte} of LPDDR5X at
    [[FILL: ${\approx}\SI{273}{\giga\byte\per\second}$]], presenting a single
    CUDA GPU to the workload. Unlike the M2~Pro, GB10 is not a single shared
    bus: the Grace CPU and Blackwell GPU are separate dies joined by a coherent
    NVLink-C2C interconnect, so it is a coherent unified-memory \emph{system}
    rather than one physical pool
    %
    ([[FILL: GPU SM count; CPU core count]]).
    %
    Software: PyTorch \texttt{2.10.0+cu130} and \texttt{torchvision
    0.25.0+cu130} (the CUDA \texttt{cu130} wheels required for the
    \texttt{sm\_121}/aarch64 target; the default wheel is CPU-only),
    \texttt{transformers~5.2}, BF16 (no quantization). We confirm GPU residency
    of every model stage with \texttt{nvidia-smi pmon}
    ([[FILL: observed per-process SM activity during generation]]).
\end{description}

\sysname{} runs the same declarative pipeline definitions on both machines,
changing only the per-stage backend and \texttt{device} fields in the
configuration; no code changes are required (Section~\ref{sec:framework}).
Because the M2~Pro runs 4-bit MLX weights while DGX~Spark runs BF16 PyTorch
weights, the two backends are \emph{not} numerically identical models. We
therefore draw cross-device conclusions only about the \emph{direction} of a
systems effect (e.g.\ whether decomposition helps), never about absolute
latency or quality across devices; absolute comparisons are reported within a
device. All model weights and datasets are pre-fetched before measured runs, so
no first-run download is timed. We monitor package power and temperature
(\texttt{macmon} on the M2~Pro; \texttt{nvidia-smi}/\texttt{dcgmi} on the NVIDIA
host) and exclude runs that show thermal throttling, so sustained-load results
are not confounded. [[FILL: confirm discard-vs-report throttling policy.]]

\paragraph{Scope.} We restrict the evaluation to these two single-node
machines. Larger discrete-memory datacenter accelerators (e.g.\ A100/H100) and
multi-node scaling are out of scope and discussed as future work
(Section~\ref{sec:conclusion}); the datacenter end of the spectrum invoked in
Section~\ref{sec:intro} is represented here by the GB10's large unified memory
rather than a discrete-memory server accelerator.
% REVIEWER (D1): reconcile this with the Intro's "high-performance data center
% nodes" framing -- consider softening the Intro to "high-capacity single-node
% platforms," or add a discrete-memory GPU.

% ---------------------------------------------------------------------
\subsection{Instrumentation and Metrics}
\label{sec:metrics}

\paragraph{Measurement pipeline.} Every run emits two artifacts. (i)~A
\emph{timing trace}: each stage and the pipeline driver log
\texttt{start}/\texttt{end} events for their \texttt{prepare} and \texttt{run}
phases as timestamped lines; per-query events additionally carry the query id,
submission timestamp, epoch, and batch index. (ii)~For the case studies, a
per-query JSONL sidecar written by \texttt{TerminalCapture}, a built-in sink
stage that records the question, golden answers, retrieved documents, and
generated answer; the instrument experiments
(\S\ref{sec:exp-noop},~\S\ref{sec:exp-modularity}) emit only the timing trace.
All durations are taken from a single monotonic clock (\texttt{time.perf\_counter};
\texttt{time.get\_clock\_info} reports a resolution of [[FILL: value]] on the
M2~Pro and \SI{1}{\nano\second} (\texttt{clock\_gettime(CLOCK\_MONOTONIC)}) on
GB10); wall-clock time (\texttt{time.time}) is
used only for cross-process trace alignment, never for per-stage latency.

In parallel, RadT attaches host-level resource listeners to each pipeline
\emph{process} (\texttt{macmon} for Apple-silicon power/thermals;
\texttt{nvidia-smi}, \texttt{dcgmi}, \texttt{iostat}, and \texttt{top} on the
NVIDIA host), and \sysname{} records MLflow trace spans whose flow identifiers
link a query's per-stage slices into a single Perfetto trace. Because isolation
is at the process boundary --- one OS process per pipeline, with the stages of
a pipeline running as threads inside it --- these listeners attribute resource
use per pipeline even when pipelines are collocated. All quantitative metrics
below are computed offline from the timing traces and JSONL sidecars by the
analysis scripts in \texttt{evaluation/scripts/}.

\paragraph{Performance metrics.} From the per-query timing intervals we report:
\begin{itemize}
  \item \textbf{Per-query latency} --- wall time from a query's pipeline
    \texttt{run}-start to its \texttt{run}-end; we report median (with the
    bootstrap CI below), mean, and max, and report p95/p99 only when the pooled
    sample across the $R$ runs reaches at least \num{100} queries,
    since tail quantiles are not estimable at smaller sizes.
  \item \textbf{Throughput} --- completed queries divided by the run's
    makespan, $\max(\text{end}) - \min(\text{start})$ over all queries. This is
    MLPerf Inference's \emph{Offline} throughput; we do not report MLPerf
    \emph{Server} (SLO-constrained sustainable QPS). The Poisson scheduler is
    used to induce stage overlap, not to certify a sustained service rate.
  \item \textbf{Stage concurrency} --- a sweep line over all stage \texttt{run}
    intervals yields the number of stages whose intervals overlap in
    wall-clock time. Because stages run as Python threads, this measures
    temporal \emph{overlap}, not CPU parallelism: pure-Python work is
    serialized by the GIL, whereas the model and retrieval stages spend their
    time inside GIL-releasing native kernels (Torch/MLX/Metal, FAISS) and so
    overlap genuinely. We therefore read concurrency together with the
    per-device occupancy below, which attributes work to hardware engines.
  \item \textbf{Per-device occupancy} --- for each device (CPU / GPU via
    \texttt{mps} / ANE on Apple; CPU / CUDA GPU on NVIDIA) we take the
    \emph{union} of the run intervals of the stages mapped to it and report the
    merged time and its fraction of makespan. We label this \emph{occupancy}
    (residency), not hardware utilization: it measures how long a device hosts
    a running stage, not how busy its execution units are. For true compute
    utilization we cross-reference \texttt{nvidia-smi}/\texttt{dcgmi}
    SM-activity counters on the NVIDIA host; the Apple Neural Engine exposes no
    per-process utilization counter, so on the M2~Pro occupancy is the only
    available per-engine signal (a stated limitation). Occupancy fractions can
    exceed 100\% in aggregate across devices when collocated stages run
    concurrently.
\end{itemize}

\paragraph{Framework-overhead metrics.} For the instrument-fidelity experiments
we additionally report per-stage overhead (mean per-query latency divided by
pipeline depth, in \si{\micro\second}), the inter-stage transition cost (the
gap between one stage's \texttt{end} and the next stage's \texttt{start}), and,
for the modularity study, the relative overhead
$(\,t_{\text{\sysname}} - t_{\text{baseline}})/t_{\text{baseline}}$ of per-step
latency against a hand-written baseline. Here the unit of work is one training
\emph{step} (one batch) driven by a single in-flight query, so ``per-step
latency'' is the per-query latency above measured at batch granularity. Warm-up
is handled per experiment and stated with each
(\S\ref{sec:exp-noop},~\S\ref{sec:exp-modularity}) rather than by a single
global rule.

\paragraph{Quality metrics (RAG case studies).} For pipelines that produce
answers we report, from the JSONL sidecar, the standard accuracy metric for each
task family: \textbf{exact match (EM)} and \textbf{token-level F1} for the
text-RAG tasks~\cite{rajpurkar2016squad,yang2018hotpotqa}, and the standard
\textbf{VQA accuracy} (soft match over the annotator answer
set)~\cite{antol2015vqa,goyal2017vqav2} for the OK-VQA pipeline. Our current
analysis scripts compute an \textbf{answer-containment rate} --- the fraction of
queries whose normalized golden answer appears as a substring of the normalized
generated answer --- which we report as a lenient secondary proxy; the EM/F1 and
VQA-accuracy scorers extend this pass over the same captured golden answers.
Containment is gameable by verbosity (a longer answer is more likely to contain
the golden string), and the topology study compares models of different
verbosity, so we report it alongside mean answer length and flag any advantage
that co-varies with length; it is a coarse regression detector, not an accuracy
score. We also report the \textbf{answered rate} (queries yielding a non-empty
answer that does not hit the retry-exhaustion sentinel) and, to confirm two
topologies attempt the same task, a \textbf{cross-pipeline parity} score: the
mean token-level Jaccard overlap of answers to shared questions (we also report
the count matched at Jaccard${\ge}0.5$). Parity measures \emph{agreement}, not
correctness: two arms may agree on the same wrong answer. Read together with
EM/F1, it flags whether a faster configuration has silently degraded answers,
but it does not certify equal quality across the 9B and 4B models of
\S\ref{sec:exp-topology}.

\paragraph{Statistical methodology.} Each configuration is run $R$ times
($R=5$) with fixed, distinct seeds controlling query order and the
Poisson arrival process. Latency distributions are pooled across the queries of
all $R$ runs; we report the median and, where the pooled sample size permits, a
nonparametric bootstrap 95\% confidence interval ($10^4$ resamples) rather than
the mean, since per-query latency is heavy-tailed. For rate metrics (answered
rate, containment rate, EM/accuracy) we report Wilson 95\% intervals on the
underlying counts and call a difference between two arms significant only when
their intervals are disjoint. Throughput is reported as the across-run median
with its CI. Unless stated otherwise, generation is greedy: we set
\texttt{do\_sample=false} (temperature~0) explicitly, independent of backend
defaults, so answers are deterministic given the model and prompt; we verify
determinism with a repeat-run identical-output check and use the same decoding
policy for every arm of a comparison. This removes sampling as a confound in the
topology and parity comparisons.

\paragraph{Reproducibility.} Models are pinned by Hugging Face commit revision
([[FILL: hashes for Qwen3.5-9B and Qwen3.5-4B, CLIP-ViT-L/14]]); the 4-bit MLX
weights are produced from those revisions with [[FILL: quantization recipe /
bit-layout]]. Datasets are pinned by version and snapshot date
([[FILL: rag-mini-wikipedia, HotpotQA/FlashRAG split, OK-VQA, COCO-Karpathy]]);
FAISS and ChromaDB indices are built with [[FILL: index type, metric, build
params]]. The Poisson load generator uses a per-experiment arrival rate, chosen
so the entry queue backs up in pipelined mode (VQA $\lambda=2.0$; Self-RAG
$\lambda=4.0$ on the CUDA host and $\lambda=1.0$ on the 4-bit MLX host),
confirmed by the measured stage-concurrency mean; serial runs set
\texttt{serialize\_queries=true}. All configurations, seeds, and analysis
scripts are released at [[FILL: URL]] so that the reported numbers are
reproducible from the released traces.

% ---------------------------------------------------------------------
\subsection{Experiments}
\label{sec:experiments}

We run four experiments. Two are instrument-fidelity checks that bound the
measurement overhead \sysname{} introduces, so that overhead reported for real
workloads can be attributed to the workload rather than the harness; the other
two --- our primary evidence --- use \sysname{} to surface end-to-end,
collocation-aware behavior that single-stream, single-engine benchmarks cannot
observe. Each is specified by its \textbf{Question}, \textbf{Design}, and
\textbf{Metrics}; results appear in Section~\ref{sec:results}.

% .....................................................................
\subsubsection{Framework overhead}
\label{sec:exp-noop}

\textbf{Question.} \sysname{} executes every stage in its own thread and moves
data between stages through queues. Before trusting any end-to-end number, we
bound the cost of that machinery: how much latency does each additional stage
add, and does passing data by reference avoid per-stage copy cost?

\textbf{Design.} We construct synthetic pipelines of pass-through (no-op)
stages that perform no computation, so all measured time is framework overhead.
Two sweeps: (i)~a \emph{depth} sweep over chain length
$N \in \{1,2,3,4,5,6,7,8,9,10,16,32,50,64,100\}$ to isolate per-stage dispatch
and queue hand-off cost; and (ii)~a \emph{payload} sweep at fixed depth~10 that
passes payloads of \SI{1}{\kibi\byte}, \SI{1}{\mebi\byte}, and
\SI{10}{\mebi\byte} under two regimes --- reference passing vs.\ a stage that
deep-copies its payload each hop --- to isolate data-movement cost.
Configurations are generated programmatically
(\texttt{noop\_chain\_generator.py}). The load generator runs \emph{closed-loop}
(\texttt{OfflineLoadScheduler}, exactly one query in flight at a time), so this
experiment measures per-stage dispatch and reference-vs-copy cost in isolation;
it deliberately does \emph{not} measure queue backlog (queueing), which is
exercised in the contention case studies below. We separate the framework's
overhead into two costs and measure both: its \emph{core dispatch} (thread
wake-up, queue hand-off, and the timing log), measured with the MLflow tracing
layer disabled (\texttt{CHOREO\_DISABLE\_TRACING}); and the \emph{tracing layer}
itself (three spans per stage per query), measured with tracing on, so its cost
can be read against real stage latencies (Section~\ref{sec:results}). All
per-stage timings use the monotonic \texttt{perf\_counter\_ns} clock emitted
alongside the wall-clock trace; because the stages of a pipeline are threads in
one process, this clock is directly comparable across stages. We discard the
first query of each run as warm-up, pool the remainder across the $R$ runs, and
report the median with the bootstrap CI of Section~\ref{sec:metrics}. Because
no accelerator kernel runs in a no-op stage, this overhead is a property of the
framework and the CPython interpreter rather than of the device; we therefore
bound it on a single representative machine (the M2~Pro), which also avoids
conflating interpreter and OS-scheduler differences with the quantity of
interest.

\textbf{Metrics.} Median per-query latency, per-stage overhead
(\si{\micro\second}), inter-stage transition cost, and the reference-vs-copy
latency ratio across payload sizes. We read these against three claims.
\emph{(i)~Depth-flatness:} per-query latency is linear in depth and the
inter-stage transition cost does not grow with depth, so per-stage overhead does
not accumulate and arbitrarily deep, non-linear graphs incur no measurement
distortion. \emph{(ii)~Zero-copy:} reference passing is constant in payload size
($O(1)$) while the deep-copy arm grows linearly ($O(\text{payload})$); the
copy arm is the counterfactual for a framework that serializes data between
stages. \emph{(iii)~Overhead-in-context:} the core-dispatch tax, and even the
heavier tracing-layer cost, are a negligible fraction of a single real model or
retrieval stage (Section~\ref{sec:results}), so overhead in the case studies is
attributable to the workload rather than the harness. \emph{Threats to validity.}
Because no-op
stages never release the GIL, the per-stage figure is a worst-case,
GIL-serialized hand-off cost and an upper bound on what real
(native-kernel-bound) stages incur; the copy arm measures in-process
\texttt{deepcopy} (memory-copy) cost and therefore lower-bounds the
serialization cost of the process-level isolation the case studies rely on.
This experiment tests whether per-stage overhead is small and flat in depth and
whether reference passing is constant in payload size; if so, overhead in real
workloads is attributable to the workload rather than the harness.

% .....................................................................
\subsubsection{Modularity overhead}
\label{sec:exp-modularity}

\textbf{Question.} Empty stages being cheap (\S\ref{sec:exp-noop}) does not
guarantee that wrapping a \emph{real} workload in \sysname{}'s graph, queue, and
thread structure is cheap. We measure whether this structure inflates runtime
relative to a monolithic, hand-written implementation.

\textbf{Design.} We express the same training workload two ways and compare
per-step latency distributions: (a)~a standalone PyTorch script with no
framework, and (b)~an equivalent \sysname{} pipeline. The workload is
transfer-learning fine-tuning of EfficientNetV2-S~\cite{tan2021efficientnetv2}
on Imagenette~\cite{imagenette} (frozen backbone, replaced classifier head,
Adam, batch size~8). The two implementations share identical data loading
(\texttt{batch\_size}=8, \texttt{shuffle}, \texttt{drop\_last}, and matched
\texttt{num\_workers}), the identically frozen backbone, and the same Adam /
cross-entropy step (the same \num{12810} trainable parameters); both run the
\texttt{train} split exclusively and both call
\texttt{torch.\{mps,cuda\}.synchronize()} at the end of each step, so logged time
reflects hardware completion rather than asynchronous kernel dispatch, leaving the
surrounding framework as the only difference. The measured per-step bracket
excludes data loading in both implementations. To isolate the wrapper rather than a
data-path execution difference, we set \texttt{num\_workers}=0 in both, so neither
prefetches a batch concurrently with the measured step (with background prefetch the
serialized pipeline and the monolith contend for memory bandwidth differently, which
confounds the per-step comparison); the \sysname{} data loader remains a separate,
independently-measured stage. Per-step latency is taken from the monotonic
\texttt{perf\_counter\_ns} clock. Each run is a single continuous epoch; we verified
that the per-step latency is \emph{flat} within a run (the device reaches a stable
clock state immediately and does not drift), so we simply discard the first
\num{200} steps as warm-up (kernel autotuning and first-call compilation), pool the
remainder across $R{=}5$ runs, and report the median with a two-sample bootstrap CI.
As in \S\ref{sec:exp-noop} we measure two arms --- the framework's \emph{core
dispatch} (MLflow tracing disabled) and the full \emph{instrumented} cost (tracing
on); the hand-written baseline carries no tracing. We \emph{interleave} the arms run
by run so they see identical conditions, and run on both DUTs (M2~Pro \texttt{mps}
and DGX~Spark \texttt{cuda}), reading overhead within a device. Because batch-8
EfficientNetV2-S is strongly GPU-bound (per-step compute $\gg$ framework
overhead), the overhead is naturally a small fraction of step time; rather than
contrive a stress configuration, we report the overhead in \emph{absolute} terms
(\si{\micro\second}) as well as a ratio, so it reads as a fixed, step-size-%
independent cost, and we rely on \S\ref{sec:exp-noop}, which establishes directly
that the per-stage framework cost is fixed and does not scale with the work a
stage does.

\textbf{Metrics.} Pooled per-step latency (median, with a nonparametric bootstrap
95\% CI) for each implementation, and the overhead of \sysname{} over the
baseline as both an \emph{absolute} cost (\si{\micro\second}) and a ratio, with a
two-independent-sample bootstrap CI. We report the absolute cost alongside the
ratio so its magnitude is explicit: a per-step overhead whose CI excludes zero is
statistically real, yet at \si{\micro\second} scale against a multi-millisecond
step it is practically negligible. \emph{Threats to validity.} A modularity
wrapper can only add work, so the overhead must be non-negative; a measured
\emph{negative} value therefore signals a confound rather than a framework
speedup, and we treat it as a bug to be found, not a result. We encountered two:
an earlier wall-clock-based reading was corrupted by differencing non-monotonic
timestamps across the pipeline's threads (removed by the monotonic
\texttt{perf\_counter}); and an earlier run schedule (all baseline runs first,
with background DataLoader prefetch) let a data-path execution difference, not the
wrapper, dominate (removed by interleaving the arms, matching \texttt{num\_workers},
and measuring at steady state). With those controlled, the overhead is a small,
correctly-signed positive. This experiment tests whether end-to-end modularity
overhead is a negligible fraction of the step it wraps.

% .....................................................................
\subsubsection{Unified-memory bandwidth contention (multimodal VQA)}
\label{sec:exp-bandwidth}

\textbf{Question.} On an Apple-silicon unified-memory SoC, the GPU, Neural
Engine, and CPU share one physical memory and its bandwidth. Single-engine
accelerator-throughput benchmarks cannot see the resulting contention. When we
map the stages of one workload across different engines, is the apparent benefit
of that heterogeneity real, or does it evaporate once queries overlap and the
engines contend for shared memory bandwidth? (This experiment runs on the
M2~Pro only; the Neural Engine has no GB10 analogue.)

\textbf{Design.} A single image-grounded VQA pipeline ---
CLIP-ViT-L/14~\cite{radford2021clip} vision encoding $\rightarrow$ FAISS
retrieval~\cite{johnson2019faiss} over \num{10000} COCO-Karpathy
captions~\cite{lin2014coco,karpathy2015deepvs} (top-3) $\rightarrow$ 4-bit MLX
Qwen3.5-9B (\texttt{mlx-community/Qwen3.5-9B-OptiQ-4bit}~\cite{qwen35}) answer
generation on the GPU via MLX --- over OK-VQA~\cite{marino2019okvqa} questions.
We run it as a $2\times2$ over \emph{accelerator mapping} and \emph{schedule}:
\begin{itemize}
  \item \emph{Mapping~A (collocated):} CLIP and the LLM both on the GPU
    (\texttt{mps}); retrieval on CPU.
  \item \emph{Mapping~B (heterogeneous):} CLIP exported to Core~ML and run on
    the Neural Engine (\texttt{ane}); the LLM on the GPU; retrieval on CPU.
  \item \emph{Schedule:} \emph{pipelined} (Poisson arrivals at
    \texttt{rate}=2.0~queries/s, above the ${\approx}3$~s/query LLM service
    time, so several queries stay in flight) vs.\ \emph{serial} (one query
    end-to-end at a time). The serial cells are produced with the load
    generator's \texttt{-{}-serialize} override; the configurations are
    otherwise identical.
\end{itemize}
\emph{Confound.} Mapping~B changes more than placement: the encoder switches
from a PyTorch/Metal model (\texttt{mps}) to a Core~ML FP16 package on the ANE,
so engine, runtime, and numerical precision co-vary, and the ANE incurs a
first-call compilation cost. A raw B-vs-A delta therefore conflates these
factors. We read the bandwidth signal primarily \emph{within} a mapping (serial
vs.\ pipelined for the same encoder), discard the ANE first-call compilation
query as warm-up and report it separately, and use cross-pipeline answer parity
(matched per question id from the JSONL sidecar) as a coarse check on downstream
divergence from the FP16 embeddings --- it does not directly measure
retrieval-set overlap. Each cell runs \texttt{max\_queries}=30; we report tail
latencies only on the distribution pooled across the $R$ runs
(\S\ref{sec:metrics}), never per run.

\textbf{Metrics.} Per-query latency, throughput, stage concurrency, and
per-device occupancy (CPU/GPU/ANE) for each cell; VQA accuracy and answer parity
to confirm the arms are comparable. The diagnostic of interest is whether the
heterogeneity advantage collapses from serial to pipelined --- behavior
\emph{consistent with} a bandwidth-bound rather than compute-bound regime. We
corroborate it with a direct DRAM-bandwidth/power read
([[FILL: \texttt{powermetrics}/\texttt{macmon}]]) where available; absent that,
we rule out the principal alternatives (non-overlapping ANE/GPU execution,
command-buffer queueing) using the per-stage concurrency and occupancy traces.

% .....................................................................
\subsubsection{Pipeline topology: monolithic vs.\ decomposed agentic RAG}
\label{sec:exp-topology}

\textbf{Question.} For an agentic, self-reflective RAG
workload~\cite{asai2024selfrag}, is it better to use one large model that
performs every sub-task (relevance grading, answering, hallucination checking)
in a single call, or several smaller specialist models composed into a graph?
Such exploratory, ablation-style restructuring is difficult to express in
fixed-topology benchmarks but is directly supported by \sysname{}'s declarative,
non-linear graphs.

\textbf{Design.} We implement a Self-RAG pipeline in two topologies that target
the same task and swap between them by changing only the configuration:
\begin{itemize}
  \item \emph{Monolith:} one large model (Qwen3.5-9B) does grade + answer +
    hallucination-check in a single structured pass; a router validates the
    output and optionally loops through a query rewriter (a back-edge in the
    graph, bounded by a retry budget).
  \item \emph{Decomposed:} three smaller specialist instances (Qwen3.5-4B;
    grader / generator / hallucination-checker, with the rewriter sharing the grader)
    split the job into separate stages connected by routers, so cheap
    grade/check calls for one query can overlap the expensive generation for
    another.
\end{itemize}
The decomposed topology exercises \sysname{}'s fan-out routing, fan-in merge
policies, and a feedback cycle (the retrieve--grade--rewrite loop), and is a
within-pipeline collocation case: three models co-resident on one accelerator,
profiled per stage. We run both topologies on two task difficulties (factoid:
\texttt{rag-mini-wikipedia}; multi-hop: HotpotQA~\cite{yang2018hotpotqa}),
retrieving from ChromaDB~\cite{chroma}, on both DUTs (DGX~Spark with BF16
HF/PyTorch models; M2~Pro with 4-bit MLX models). Retrieval is configured per
difficulty and held fixed across both topologies: factoid retrieves top-3 over
3{,}200 \texttt{rag-mini-wikipedia} passages, while multi-hop retrieves top-5
over a 20{,}000-passage corpus built from the gold and distractor contexts of
the first 300 HotpotQA-dev questions (via the FlashRAG redistribution). Because
retrieval depth and corpus size differ across difficulties, the
monolith-vs-decomposed contrast is read within a difficulty, not across.

\emph{Confound and control.} The monolith uses one 9B model while the
decomposed graph uses three 4B instances, so topology, model \emph{size}, and
\emph{call count} (one long generation vs.\ three short, overlappable ones)
co-vary. To attribute an effect to topology rather than to using a smaller
model, we add a size-matched control --- a 4B monolith issuing the three calls
sequentially --- so the decomposed-vs-control contrast isolates
overlap/specialization at fixed parameter size [[FILL: confirm control arm
run]]. Load generation is \emph{not} identical across DUTs: the CUDA host uses
\texttt{rate}=4.0 with 30~queries, while the slower 4-bit MLX host uses
\texttt{rate}=1.0 with 10~queries, each chosen so the entry queue backs up in
pipelined mode. We therefore compare topologies \emph{within} a DUT and treat
cross-DUT numbers as portability checks, not head-to-head latency. Serial cells
use the \texttt{-{}-serialize} override.

\textbf{Metrics.} The full performance set (per-query latency, throughput, stage
concurrency, per-device occupancy, LLM calls per query, resident model memory)
and the quality set (EM/F1, answered rate, answer parity). Because the arms use
different model sizes, we condition every speed comparison on the EM/F1 of the
two arms: a speed result obtained at degraded accuracy is reported as a
trade-off, not as a free win. The contribution here is methodological as much as
empirical: a single declarative reconfiguration, profiled end-to-end and
per-stage across two machines, yields a trade-off study that an isolated,
single-graph benchmark cannot produce.

\paragraph{On collocation.} We do not run a separate independent-workload
interference experiment in this paper. Collocation is instead exercised
\emph{within} the case studies --- the VQA pipeline collocates CLIP and the LLM
on one GPU (\S\ref{sec:exp-bandwidth}), and the decomposed Self-RAG pipeline
collocates three model instances on one accelerator
(\S\ref{sec:exp-topology}) --- and \sysname{}'s per-process, per-stage
instrumentation is what makes the resulting contention measurable. A controlled
study of independent collocated workloads (e.g.\ foreground inference contending
with a background fine-tune) is a natural next use of the framework and is left
to future work.
% REVIEWER (D1, HIGH): the Intro headlines "collocation-aware profiling" and
% "launching an arbitrary number of concurrently running pipelines" as a core
% contribution; declining the inter-pipeline experiment here reads as the main
% claim being unsupported. Strongest fix: add ONE cell launching a Poisson VQA
% inference pipeline collocated with the offline EfficientNet fine-tune
% (torchvision_mixed.yml exists), reporting foreground p95 degradation vs.
% isolation and per-PROCESS occupancy. Otherwise soften Contribution 4 and the
% framework prose. Also unaddressed: no first-class scalability sweep (Intro
% claims "no-code scalability"/"stress hardware of any size" -- a model-size
% sweep via the declarative `model` field would substantiate it), and no
% baseline-tool (MLPerf) end-to-end-vs-isolated contrast despite the paper's
% "MLPerf falls short" thesis (pipeline_configs/mlperf/ configs exist).

% =====================================================================
%  REQUIRED CITATIONS (add to .bib; keys used above are placeholders):
%   rajpurkar2016squad  - SQuAD (EM/F1)                 - Rajpurkar+ EMNLP'16
%   yang2018hotpotqa    - HotpotQA                       - Yang+ EMNLP'18
%   antol2015vqa        - VQA accuracy metric (OK-VQA)    - Antol+ ICCV'15
%   goyal2017vqav2      - VQA v2 accuracy                 - Goyal+ CVPR'17
%   tan2021efficientnetv2 - EfficientNetV2               - Tan & Le ICML'21
%   imagenette          - Imagenette dataset             - Howard (fast.ai)
%   radford2021clip     - CLIP                           - Radford+ ICML'21
%   johnson2019faiss    - FAISS                          - Johnson+ IEEE BigData'19
%   lin2014coco         - MS COCO                        - Lin+ ECCV'14
%   karpathy2015deepvs  - Karpathy caption split         - Karpathy & Fei-Fei CVPR'15
%   marino2019okvqa     - OK-VQA                         - Marino+ CVPR'19
%   qwen35              - Qwen3.5 + OptiQ-4bit MLX card
%   asai2024selfrag     - Self-RAG                        - Asai+ ICLR'24
%   chroma              - ChromaDB (project/docs)
%  ALSO cite where referenced elsewhere in the paper:
%   MLPerf Inference (Reddi+ ISCA'20; Offline/Server scenarios), MLPerf Tiny
%   (Banbury+ MLSys'21), TPCx-AI (Brücke+ VLDB'23), AIBench (Gao+), DAWNBench
%   (Coleman+ '17), MLX (Apple ml-explore), Core ML/ANE (Apple docs), Perfetto,
%   MLflow (Zaharia+ '18), Phi-3-mini (Abdin+ '24), RAG (Lewis+ '20),
%   DGX Spark/GB10 datasheet (128GB LPDDR5X, 273 GB/s, NVLink-C2C), and the
%   draft's own margin TODOs: AISBench + the HPI TPCx-AI VLDB preprint.
% =====================================================================
